\section{Introduction}

Flaky tests — tests that pass and fail nondeterministically on the same code
version — are among the most corrosive reliability problems in modern software
engineering. Surveys of industrial projects estimate that roughly one test in
ten exhibits flakiness~\cite{luo2014empirical,lam2019understanding}, and flaky
failures are consistently reported as a leading cause of developer distrust in
CI and of delayed releases~\cite{google2023flaky,seka2023survey}. A developer
can fix a flaky test only after two things happen: a failure is
\emph{reproduced} (ideally locally, in an environment the developer controls),
and the \emph{root cause} is understood (which state, order, timing condition,
or nondeterministic behavior produces it). Both steps are expensive.
Flakiness manifests probabilistically, under conditions that are hard to guess
(order of test execution, wall-clock timing, machine load, locale, hash
ordering), and a developer who clicks ``re-run'' in CI and sees green is left
with nothing to debug.

The research community has built increasingly capable tools for this problem,
but each addresses a slice of it. Classification tools predict
\emph{whether} a test is flaky~\cite{king2023bayesian,pinto2020understanding,
kim2024large}. Specialized reproducers target a single flakiness
category: order-dependent failures~\cite{lam2019idflakies,parry2022survey},
timing-dependent failures~\cite{chen2024tdrepro}, or non-idempotent
outcomes~\cite{shen2024nidebugger}, and typically require the failing
test's log or stack trace as input to imitate~\cite{chen2024tdrepro}. The
newest generation of LLM-based debugging assistants~\cite{lewis2020retrieval,
zhang2024llamadebug} is designed for deterministic bugs, where the failure is
reproducible on demand; flaky failures break this assumption because the
assistant may never observe the bug it is asked to fix.

In this paper we argue that flaky-test reproduction is a natural fit for
\emph{closed-loop LLM agents} — systems that propose an action, execute it in
the real environment, observe the result, and adapt — and we introduce
\textbf{FlakeAgent}, the first category-agnostic agent for reproducing and
diagnosing flaky test failures from cold start. FlakeAgent receives only a
test identifier and access to the project source; it receives
\emph{no} failure log, no stack trace, and no category hint. In a bounded
loop (at most six LLM calls and eight local executions per test, under a
45-minute wall-clock cap), it iteratively forms a causal hypothesis about the
flakiness mechanism, selects one of twelve executable strategies
(order permutations, same-JVM double runs, sleep injection, CPU/IO stress,
heap limits, environment manipulation, parallel execution, network blocking,
and NonDex-based nondeterminism injection), executes it locally, and observes
the structured result (exit status, failed tests, assertion messages, stack
traces, durations). When a failure is reproduced, FlakeAgent switches to a
diagnosis step that emits (i) a flakiness category
(\emph{order-dependent, time-dependent, non-idempotent-outcome,
implementation-dependent}), (ii) a root-cause hypothesis in natural language,
and (iii) file--line \emph{evidence} with code snippets drawn from the source
and the observed failure — evidence we evaluate against developer fix commits.

We evaluate FlakeAgent on a benchmark drawn from
\emph{ReproFlake}~\cite{zhang2026reproflake}, a recent dataset of 1{,}115
machine-reproducible flaky tests in real Java projects, shipped with the
official reproduction mechanisms for each category and with the developer fix
commit for each test. We select 116 tests, balancing all four categories (30
order-dependent, 36 time-dependent, 25 non-idempotent-outcome, and 25
implementation-dependent tests across 46 project snapshots); a validation gate
confirms the dataset's official mechanism reproduces each test in our
environment, yielding 92 gate-passed tests (20 order-dependent, 31
time-dependent, 25 non-idempotent-outcome, 16 implementation-dependent) on
which all methods are evaluated. We compare FlakeAgent against five baselines sharing the
same execution budget: naive rerun, uniform random strategy sampling, random
suite-order shuffling, a deterministic sleep-injection scan, and a single-shot
LLM prompt without a loop. We additionally ablate the agent's reflection loop,
the breadth of its strategy space, and its source-code grounding, and we
report the first cross-model comparison of flaky-test reproduction across six
2026-era frontier LLMs accessed via vendor APIs.

Our results show that FlakeAgent reproduces 47.8\% of benchmark failures
within budget (20.0\% of order-dependent, 12.9\% of time-dependent,
100.0\% of non-idempotent, 68.8\% of implementation-dependent),
outperforming the strongest non-agent baseline by 20.6 percentage points,
and that its evidence-grounded diagnoses cite file--line pairs touched by the
developer fix in 36.4\% of cases. Ablations attribute a large share of
the gain to the closed loop rather than to the underlying model alone, and
the cost analysis shows reproduction at a median of 0.006 USD per test —
two orders of magnitude cheaper than the brute-force rerun budget the same
tests required in the benchmark's own construction.

This paper makes the following contributions:
\begin{itemize}
\item \textbf{FlakeAgent}, the first category-agnostic closed-loop LLM agent
for flaky-test reproduction and evidence-grounded root-cause diagnosis from
cold start (no failure logs), with a bounded call/execution/wall-clock budget;
\item \textbf{Twelve executable reproduction strategies} spanning order
manipulation, timing perturbation, resource pressure, environment
manipulation, and nondeterminism injection, implemented executor-side so that
the model only selects parameters;
\item \textbf{A benchmark and evaluation methodology} built on ReproFlake —
116 selected tests (92 gate-passed) across four flakiness categories with a
reproducibility gate, five shared-budget baselines, four ablations, and a
six-model comparison on 2026-era frontier LLMs;
\item \textbf{Evidence-grounded diagnosis evaluation}: category accuracy and
file--line evidence checked against developer fix commits, including a
cost breakdown of agentic reproduction per test and per failure.
\end{itemize}
